\section{立案}
\begin{enumerate}
  \item 启动条件。
        \begin{enumerate}
          \item 公诉（公、检）。有犯罪事实 + 需要追究刑事责任。
          \item 自诉。亲告；法院管辖；被害人告诉；明确的被告人、具体的诉讼请求、证据。\zd{不同于民诉}。
        \end{enumerate}
  \item 当事人。报案、控告、举报、自首。
        \begin{enumerate}
          \item \bs{报案}。是任何人都可以来报，\zd{对案不对人}。
          \item \bs{控告}。针对具体的人，\zd{被害人自已}去指控。
          \item \bs{举报}。简称案外人来举报。针对具体的人，\zd{被害人以外}（不限于人，单位也可）去指控。
          \item 控告人救济。\zd{不同的角色救济途径的设计不一样}，法律\zd{只给控告人}准备了很多救济途径。
                \begin{enumerate}
                  \item \zd{找公}安复议 \yjt 上一级复核。
                  \item \zd{找同级检}察院申诉 \yjt 应当受理，并进行审查。
                  \item \zd{找法}院。公转自。
                \end{enumerate}
        \end{enumerate}
  \item 公、检。
        \begin{enumerate}
          \item 材料都应当接受。
                \begin{enumerate}
                  \item 接受材料 \bdy 立案 \yjt 管辖判断 \yjt 无权管辖 \yjt 移送。
                  \item 情况紧急的还要先采取紧急措施。
                  \item 说明诬告陷害相关的法律责任。
                  \item 保障报案控告举报人、扭送人以及近亲属的人身安全。
                \end{enumerate}
          \item 初查 = 初步调查核实 \bdy 侦查。\zd{温柔的手段}。不限制人身、财产权利的任意性措施。
          \item 书面决定立案 \hz 不立案 \yjt 告知原因。
          \item \zd{撤案}。
                \begin{enumerate}
                  \item 一般撤案。显著轻，过时效，特赦告诉和死掉，\bs{还有无罪均撤案}。
                  \item 特殊撤案。\bs{如实供、大立功、最高检核}。其实就是认罪认罚免于刑事处罚的条件。
                  \item 终止侦查。\bs{并非他所为、共犯不处罚}。
                \end{enumerate}
        \end{enumerate}
  \item 公安。对检察院不服复议、复核。
  \item 检察院。
        \begin{enumerate}
          \item 该立/不该立。\bs{先要求说理，后通知立/撤案}。\zd{7} 日内说理、\zd{15} 日内必须/立即撤案、检察长决定。
          \item 就不立，死不悔改\zd{纠正违法通知书}。
          \item 还是不立，\zd{报上一级检察院协商同级公安去处理}。
          \item 继续迟迟不作为，\zd{立案监督案件催办函}。
          \item 后面法条没有明确的规定了。
        \end{enumerate}
\end{enumerate}
